\chapter{Thực hành đạo đức với dữ liệu}
\label{ch:M2L4}

% Lesson: Ethical Practices with Data

\section*{Lesson Notes}

THỰC HÀNH ĐẠO ĐỨC VỚI DỮ LIỆU

{\small
\begin{longtable}{>{\raggedright\arraybackslash}p{0.292\linewidth}>{\raggedright\arraybackslash}p{0.648\linewidth}}
\toprule
\textbf{Thời gian đọc (Reading Time)} & 20 phút \\
\addlinespace[2pt]
\textbf{Kiến thức trước (Prior Knowledge)} & Mô-đun 1 môn Dữ liệu Tài chính (FD), Các loại Dữ liệu, Thuộc tính của Dữ liệu, Thực hành Tốt với Dữ liệu \\
\addlinespace[2pt]
\textbf{Từ khóa (Keywords)} & Flash Crash, Trách nhiệm ủy thác (Fiduciary), Làm mượt Dữ liệu (Smoothing Data) \\
\addlinespace[2pt]
\bottomrule
\end{longtable}
}

\emph{Không phải tất cả dữ liệu đều được tạo ra một cách có đạo đức, thẳng thắn, đầy đủ và minh bạch. Giống như việc chúng ta cần áp dụng tư duy phản biện để đánh giá chất lượng, độ chính xác và tính đầy đủ của dữ liệu, chúng ta cũng nên sử dụng các thực hành đạo đức tốt khi cung cấp dữ liệu cho người khác hoặc sử dụng nó trong các mô hình của riêng mình. Bài học này sẽ xem xét một tá các thực hành phi đạo đức liên quan đến dữ liệu.}

Đạo đức là một phần quan trọng của tài chính. Các công ty thích đưa đạo đức vào việc đào tạo nhân viên của họ. Thật vậy, một trong những công ty kế toán lớn nhất thế giới đã tổ chức một bài kiểm tra đạo đức cho nhân viên của mình. Hóa ra là các nhân viên đã gian lận trong bài kiểm tra đạo đức của họ! Bản thân công ty đã không tiết lộ thông tin này cho cơ quan quản lý (regulatory board) (GaNun).

Bài học này đề cập đến các vấn đề đạo đức khi chúng áp dụng vào việc xử lý dữ liệu. Mỗi tiểu mục mô tả một trong những cách phi đạo đức mà mọi người xử lý dữ liệu không phù hợp: đầu tư (investing) , phóng đại, bóp méo, thu thập không phù hợp, che giấu, hiểu sai, v.v. Mỗi tiểu mục bao gồm ít nhất một ví dụ thực tế về hành vi được thực hiện bởi một cá nhân, công ty, hoặc thậm chí là một thị trường.

\section{Cố ý tạo ra dữ liệu gây hiểu lầm (Intentionally Creating Misleading Data)}

Một ngày nọ, một nghệ sĩ tên là Simon Weckert đã thu thập 99 chiếc điện thoại thông minh (Barret). Anh ta đặt chúng vào một chiếc xe kéo mà sau đó anh ta tiếp tục kéo đi bộ dọc theo các con phố của Munich, trớ trêu thay lại là con phố gần trụ sở Google ở Munich---anh ta đang hack ứng dụng Google Maps! Khoảng một giờ sau khi bắt đầu đi bộ, Google Maps đã hiển thị tình trạng giao thông ùn tắc, kéo dài---mặc dù không có bất kỳ chiếc ô tô nào trong khu vực. Nghệ sĩ này tuyên bố rằng anh ta đang chỉ ra sự đơn giản đằng sau các hệ thống này. Thực sự không có tác hại nào được gây ra trong trò chơi khăm thông minh này. Điều này khiến chúng ta đặt câu hỏi về các giả định của mình khi nhận được dữ liệu, chẳng hạn như điều kiện giao thông.

Nếu trò này được thực hiện trên các thị trường tài chính, đó sẽ là một câu chuyện hoàn toàn khác. Việc cố ý tạo ra dữ liệu gây hiểu lầm là phi đạo đức. Giống như việc chúng ta có thể tự động tuân theo các điều kiện trên một ứng dụng giao thông, một cố vấn robot (RoboAdvisor) có thể tự động tuân theo các điều kiện trên thị trường (D'Acunto và Rossi). Chúng có thể phản ứng mà không đặt câu hỏi về dữ liệu vì chúng tập trung vào tốc độ và tính hiệu quả thay vì xác thực kết quả.

Một cách trắng trợn để cố ý gây hiểu lầm cho một thị trường là giả mạo (spoof) nó. Một người như vậy đã làm điều đó vào năm 2010 là Navinder Singh Sarao. Ông ta bị quy trách nhiệm cho việc khơi mào một đợt sụt giảm mạnh của thị trường (market downturn) được gọi là Flash Crash (sụp đổ chớp nhoáng) năm 2010. Sarao đã nhận tội giả mạo thị trường và phải trả giá bằng một bản án tù cùng các khoản tiền phạt.

Giả mạo (Spoofing) là gì? Giả mạo là một kiểu lừa gạt cực đoan. Việc giả mạo xảy ra khi một nhà giao dịch đặt một lệnh lớn mà anh ta không có ý định khớp lệnh. Thay vào đó, ý định là thao túng thị trường có lợi cho mình. Ví dụ, một kẻ giả mạo (spoofer) đặt một lệnh mua lớn sẽ đặt một lệnh chào mua khổng lồ thấp hơn nhiều so với báo giá chào mua hiện tại của thị trường. Khi các nhà tạo lập thị trường và các nhà giao dịch khác, bao gồm cả các nhà giao dịch thuật toán, cảm nhận được một người mua có động lực, họ sẽ rút các giá chào bán của mình và tăng chúng lên. Những đợt tăng giá này kích hoạt các nhà giao dịch khác làm điều tương tự. Tuy nhiên, đây mới là động cơ của kẻ giả mạo từ đầu---một sự thao túng giá theo hướng mong muốn.

Trong trường hợp của ông Sarao, ông ta đã giả mạo thị trường bằng cách đặt một lệnh bán lớn (cao hơn giá chào bán) cho hợp đồng tương lai E-mini S\&P. Lệnh của ông ta đã khiến các thuật toán hạ giá của chúng, điều này đến lượt nó lại khiến các nhà giao dịch và thuật toán khác hạ giá của họ, và cứ tiếp tục như vậy. Giá cả lao dốc. Các hành động và phản ứng đã gây ra Flash Crash, một đợt sụp đổ thị trường trong ngày đã quét sạch hơn \$1 nghìn tỷ giá trị thị trường. Các thị trường khác không cho thấy sự sụt giảm, và vì vậy thị trường đã học được rằng đây là một tai nạn xảy ra một lần. Trong vòng một giờ, hầu hết thị trường đã điều chỉnh.

Việc tạo ra dữ liệu gây hiểu lầm là phi đạo đức. Các thị trường giao dịch trong thời gian thực và có thể bị gây hiểu lầm bởi những kẻ giả mạo. May mắn thay, quy định và các kiểm soát lớn hơn tồn tại ở nhiều sàn giao dịch ngày nay giúp giảm thiểu sự tái diễn của việc một nhà giao dịch giả mạo thị trường và gây ra một đợt Flash Crash khác. Việc giả mạo thị trường vẫn là một hành vi phạm tội hình sự.

\section{Lan truyền đánh giá giả (Propagating Fake Reviews)}

Lần tới khi bạn ở London, bạn có thể muốn dùng bữa tại nơi từng là nhà hàng được đánh giá cao nhất trên tripadvisor.com: The Shed (Butler). Với hình ảnh các món ăn, một khu vườn bên ngoài, và hàng chục đánh giá, The Shed dường như là một nơi hợp thời trang để ăn uống. Vấn đề duy nhất là nhà hàng này không hề tồn tại. Các đánh giá ban đầu đều là giả mạo. Trớ trêu thay, The Shed lại là nơi mọi người đổ xô đến. Những cuộc gọi điện thoại đến nhà hàng chỉ được kết nối tới máy trả lời tự động. Khi mọi người nhận ra nhà hàng này hợp thời trang đến mức nào, họ đã viết đánh giá mặc dù họ chưa bao giờ thực sự đến đó. Một phần của động lực lan truyền (going viral) là việc thuê những người nổi tiếng tạo dáng với các món ăn (vốn thậm chí không phải từ nhà hàng) để tạo ra bầu không khí uy tín và phong cách.

Đọc các đánh giá không giống như việc thực hiện thẩm định chuyên sâu (due diligence). Khi nhiều chi tiết hơn của vụ bê bối FTX được công khai, một số người nổi tiếng đang đứng ra thừa nhận rằng họ là nhà đầu tư. Nhiều người quen thuộc với ngôi sao điện ảnh Kevin Bacon. Anh ấy là một nhà đầu tư vào FTX và tuyên bố rằng anh ấy đã mất hàng triệu đô la. Ngay cả những người nổi tiếng cũng có thể trở thành nạn nhân của hiệu ứng hào quang (halo effect) hoặc sự viện dẫn thẩm quyền (appeal to authority). Các nhà tiếp thị biết sức mạnh của sự bảo chứng của người nổi tiếng, đến mức nhiều vận động viên và nhạc sĩ kiếm được từ việc bảo chứng sản phẩm nhiều hơn từ tiền lương hàng triệu đô la.

Bạn có tự hỏi có bao nhiêu đánh giá nhà hàng là giả mạo? Có những bên với lợi ích nhóm bị cám dỗ để đăng hết đánh giá rực rỡ này đến đánh giá rực rỡ khác. Một cú nhấp chuột nhanh vào người đánh giá đó cho thấy họ không đánh giá điều gì khác. Điều tương tự cũng có thể nói đối với các diễn đàn đánh giá cổ phiếu. Một số người tham gia các diễn đàn phổ biến này để chào mời cổ phiếu (hoặc để vùi dập chúng nếu họ có các vị thế bán khống). Đạo đức tốt sẽ tránh sự xuyên tạc.

\section{Không tiết lộ phí hoặc xung đột lợi ích (Failing to Disclose Fees or Conflicts of Interest)}

Theo một cuộc khảo sát dữ liệu của công ty nghiên cứu dữ liệu Morning Consult, chỉ 36\% người lớn tin tưởng các nhà quản lý đầu tư và quản lý tài sản. Phần còn lại hoặc không tin tưởng họ hoặc không có ý kiến về họ. Một phần của vấn đề là các cố vấn tài chính không tiết lộ các khoản phí hoặc tiết lộ các xung đột lợi ích. Ví dụ, một cố vấn có thể đề xuất quỹ tương hỗ A thay vì quỹ B cho một khách hàng. Mặc dù B có thể phù hợp hơn cho hồ sơ đầu tư và khả năng chịu đựng rủi ro, nhưng cố vấn lại được trả một khoản hoa hồng tốt hơn nhiều cho việc bán quỹ tương hỗ A. Việc không tiết lộ các khoản phí hoặc tiết lộ xung đột lợi ích vi phạm lòng tin ủy thác (fiduciary trust).

Trong thế giới đầu tư, đây là những hành vi phạm tội. Ví dụ, hai cố vấn robot (robo-advisor), một trong số đó có hơn \$11 tỷ tiền của khách hàng, đã đưa ra các tuyên bố sai lệch về một chiến lược liên quan đến việc thu hoạch lỗ tính thuế (tax-loss harvesting) (U.S. Securities and Exchange Commission, ``SEC Charges Two Robo-Advisers''). Những tiết lộ sai lệch này đã khiến các nhà đầu tư lầm tưởng rằng các quỹ có hiệu suất vượt trội (lợi nhuận cao hơn, rủi ro ít hơn, phí nhỏ hơn) và tự nhiên khiến chúng có vẻ hấp dẫn hơn so với các quỹ khác. Những lần khác, các nhà quản lý không tiết lộ các khoản phí (U.S. Securities and Exchange Commission, ``SEC Charges Investment Adviser'').

Đạo đức tốt là công khai mọi dữ liệu liên quan đến đầu tư: kết quả lịch sử, các khoản phí, và rủi ro. Quy định cũng giúp ích. Một số ngành bảo vệ người tiêu dùng/nhà đầu tư tốt hơn những ngành khác.

\section{Dùng dữ liệu mô phỏng thay cho dữ liệu thực mà không công khai rõ ràng (Using Simulated Data Instead of Real Data without Explicit Disclosure)}

Một cách khác để nói dối với dữ liệu là sử dụng dữ liệu mô phỏng trong khi cố tình đánh lừa rằng đó là dữ liệu thực. Ví dụ, một quỹ phòng hộ đã huy động được hàng triệu đô la bằng cách tuyên bố sai sự thật rằng họ có một thành tích (track record) dài hạn dựa trên các giao dịch thực tế sử dụng tiền thật (U.S. Securities and Exchange Commission, ``SEC Charges Father-and-Son Hedge Fund Managers''). Trên thực tế, các kết quả là từ các mô phỏng giả định được kiểm thử ngược (back-tested). Các mô phỏng có những chi tiết có thể làm sai lệch một nghiên cứu theo cách này hay cách khác. Ví dụ, một mô phỏng của một quá trình kiểm thử ngược có thể bỏ qua khoảng chênh lệch chào mua - chào bán (bid-ask spread) và tác động thị trường (market impact). Việc xuyên tạc loại dữ liệu là bất hợp pháp, và các nhà quản lý đã đồng ý trả hàng triệu đô la để giải quyết vụ án gian lận. Bản thân việc sử dụng lợi nhuận mô phỏng không phải là bất hợp pháp hay phi đạo đức. Khi sử dụng dữ liệu mô phỏng, bạn nên chỉ ra rõ ràng rằng các kết quả không phải là thực, mà chỉ là sản phẩm của các kiểm thử ngược hoặc các kịch bản. Ngay cả khi một quỹ có một thành tích hoạt động kéo dài nhiều năm, một tuyên bố miễn trừ trách nhiệm (disclaimer) được thêm vào rằng hiệu suất trong quá khứ không đảm bảo cho các kết quả trong tương lai. Gây hiểu lầm là nói dối.

\section{Thêm dữ liệu tốt vào dữ liệu xấu (Adding Good Data to Bad Data)}

Việc học sinh gian lận thì đã được biết đến nhiều, nhưng cũng có một thứ gọi là giáo viên gian lận. Câu chuyện này đã được mô tả trong cuốn sách \emph{Freakonomics: A Rogue Economist Explores the Hidden Side of Everything} của Levitt và Dubner. Để làm cho việc giảng dạy của chính họ có vẻ hiệu quả hơn, một số giáo viên ở Chicago đã gian lận bằng cách thêm điểm vào bài kiểm tra của học sinh. Khi học sinh nộp bài kiểm tra, các giáo viên sẽ điền chính xác vài câu hỏi cuối cùng, vì hầu hết học sinh không cố gắng trả lời các câu hỏi khó do thiếu thời gian hoặc độ khó. Vì vậy, các giáo viên này có thể lấy bài thi của một học sinh và dành vài giây để cải thiện điểm số. Các giáo viên không dành bất kỳ thời gian nào để xóa các câu trả lời, nhưng đã dành thời gian để điền vào một dãy đáp án trắc nghiệm đã thuộc lòng (memorized pattern) cho vài câu hỏi cuối cùng, mà hầu hết học sinh chưa bao giờ cố gắng làm. Nhìn chung, việc giáo viên gian lận dường như đã hiệu quả: điểm số của học sinh đã được cải thiện! Nhưng có một số thống kê không nhất quán về hiệu suất đã gióng lên những hồi chuông cảnh báo. Nhiều học sinh trả lời sai các câu hỏi dễ và trung bình nhưng lại làm đúng một chuỗi các câu hỏi khó. Hơn nữa, gần như tất cả các học sinh trong một lớp học đều làm đúng các câu hỏi khó giống hệt nhau. Những sự trùng hợp thật đáng kinh ngạc. Các cuộc điều tra sâu hơn cho thấy có một động lực cho các giáo viên: điểm kiểm tra được cải thiện đồng nghĩa với việc có nhiều tiền hơn cho giáo viên.

Trong ngành ngân hàng, cái giá phải trả cao hơn, và các phương pháp tinh vi hơn nhiều. Định luật Benford (Benford's Law) được sử dụng để phát hiện sự thao túng dữ liệu kế toán trong ngành ngân hàng. Định luật Benford dự đoán các tần số tương đối của các chữ số. Vui lòng đọc bài báo của Grammatikos và Papanikolaou để tìm hiểu cách Định luật Benford được áp dụng để phát hiện sự thao túng dữ liệu kế toán trong ngành ngân hàng. Người ta có thể nghĩ rằng một tỷ lệ nhỏ dữ liệu tốt sẽ không bị chú ý. Thực hành đạo đức tốt là không can thiệp vào chất lượng dữ liệu.

\section{Che giấu dữ liệu (Hiding Data)}

Thay vì thêm dữ liệu tốt vào dữ liệu xấu, tại sao không che giấu hoàn toàn một số dữ liệu? Đó là những gì đã được thực hiện trong trường hợp của một số khoản nợ công (public debt). Horn mô tả việc báo cáo dưới mức có hệ thống (systematic underreporting) của nợ công. Nghiên cứu bao gồm dữ liệu từ 50 năm. Trong các nền kinh tế mạnh mẽ, việc che giấu nợ rất dễ dàng. Tuy nhiên, trong những thời điểm kinh tế tồi tệ, các khoản nợ bị che giấu có khả năng bị phơi bày. Việc che giấu nợ có thể làm cho các chính trị gia, các kiểm soát viên, hoặc ngân sách trông có vẻ tốt hơn thực tế, và các chính trị gia có thể đưa ra những lời hứa hẹn có thể giúp họ đắc cử nhưng sẽ không thực hiện được khi các khoản thu và chi thực sự được tính toán.

Những ý định tốt có thể khiến mọi người che giấu dữ liệu. Nếu bạn đang loại bỏ các giá trị ngoại lai (outliers) khỏi một phân phối, điều đó có vẻ như đang giảm nhiễu (noise), cải thiện kết quả tổng thể. Tuy nhiên, những gì một người coi là một giá trị ngoại lai, người khác có thể coi là một trường hợp cực đoan quan trọng. Tương tự, vào một ``ngày giao dịch mỏng'' (một ngày với rất ít thanh khoản), một người có thể coi chênh lệch giá chào mua - chào bán là không đại diện cho thanh khoản tổng thể. Tuy nhiên, các giai đoạn kém thanh khoản này cần được bao gồm trong các mô hình đo lường thanh khoản và rủi ro tổng thể.

Những lần khác, dữ liệu có thể bị thiên lệch và làm sai lệch kết quả. Giả sử một nhà hàng chỉ cung cấp một biểu mẫu phản hồi cho các khách hàng quay lại. Có lẽ, một khách hàng quay lại đã thích trải nghiệm của họ và quay lại cơ sở kinh doanh đó. Tuy nhiên, còn một khách hàng chỉ đến một lần mà không có trải nghiệm tốt thì sao? Họ bị từ chối cơ hội nhận được một biểu mẫu phản hồi ngay từ đầu. Tất nhiên, nhiều trang web trực tuyến (chẳng hạn như Yelp) cung cấp cơ hội cho những khách hàng không hài lòng bày tỏ quan điểm của họ.

\section{Làm mượt dữ liệu (Smoothing Data)}

Việc làm mượt (smoothing) nghe có vẻ là một điều tốt. Trong Kinh tế lượng, bạn sẽ làm mượt các chuỗi thời gian (time series) để lọc bỏ nhiễu. Tuy nhiên, việc làm mượt lãi và lỗ (P and L) làm cho dữ liệu có vẻ ít rủi ro hơn thực tế. Việc làm mượt báo cáo dưới mức (under-reports) biến động. Báo cáo dưới mức biến động là một hiện tượng phổ biến tại các quỹ thay thế hoặc các quỹ giao dịch các chứng khoán không có định giá theo giá thị trường (mark-to-market) rõ ràng, thanh khoản (Le Marois). Phải thừa nhận rằng, nếu một chứng khoán đã không được giao dịch trong nhiều tháng hoặc khó định giá do thiếu khả năng thay thế được (fungibility), thì rất khó để đưa ra giá đánh dấu (marks) chính xác. Nếu biến động bị khai thấp (understated), thì tỷ lệ Sharpe (Sharpe ratios) bị phóng đại (overstated). Ví dụ, giả sử một chiến lược kiếm được cho người quản lý quỹ các mức lợi nhuận hàng ngày sau: +5, -1, +1, -3. Người quản lý có thể biến mức +5 thành +1, giữ lại một khoản dự trữ là 4. Khoản dự trữ đó có thể làm phẳng hoàn toàn mức -3 và giảm nhẹ mức -2 thành -1, vì vậy thay vì có những biến động tăng 5 và giảm 3, chiến lược báo cáo +1, 0, +1, 0. Điều này trông rất tốt: những chiến thắng nhỏ và không có thua lỗ. Khi được thực hiện ở quy mô lớn, việc giữ lại các khoản lợi nhuận để bù đắp các khoản lỗ sẽ làm giảm biến động của một chiến lược. Đạo đức tốt tránh việc làm mượt dữ liệu.

\section{Ngoại suy dữ liệu (Extrapolating Data)}

Việc ngoại suy (extrapolating) dữ liệu là một vấn đề. Một số mô hình có các vùng tuyến tính và sau đó nhanh chóng thay đổi thành phi tuyến tính (non-linear). Hãy nghĩ đến tác động thị trường. Trước đó trong khóa học này, chúng ta đã thấy vấn đề liên quan đến thảm họa Tàu con thoi Challenger. Vui lòng tham khảo Hình đã chuyển vị (transposed) cho Số lượng Sự cố trong bài báo của Arliss, biểu đồ biểu thị sự thất bại so với nhiệt độ. Đáng buồn thay, phân tích cho thấy rằng việc giảm 20 độ nhiệt độ (phản ánh các điều kiện nhiệt độ vào ngày phóng) làm tăng tỷ số cược (odds) thất bại lên 15 lần. Thực hành đạo đức tốt là đảm bảo các kết quả dự đoán không nằm ngoài vùng thử nghiệm.

Nguy cơ của việc ngoại suy là giả định rằng xu hướng sẽ tiếp tục bên ngoài vùng của các biến dự báo nơi nó đã được thử nghiệm. Nếu một hệ thống là tuyến tính, thì việc ngoại suy sẽ không tạo ra các vấn đề. Thách thức của tính phi tuyến tính khiến việc giả định các xu hướng tiếp tục bên ngoài vùng kiểm thử trở nên nguy hiểm.

\section{Đánh cắp dữ liệu (Stealing Data)}

Các công ty bị đánh cắp dữ liệu. Một số dữ liệu này có thể chứa thông tin khách hàng quan trọng, được sử dụng để hack vào tài khoản của họ, bòn rút tiền của họ hoặc thiết lập danh tính mới dưới tên của họ. Trong các trường hợp khác, nó có bản quyền, được cấp bằng sáng chế, hoặc được đăng ký nhãn hiệu. Trong các bài học trước, chúng ta đã thấy những nhân viên đã lấy dữ liệu (và code) từ các công ty. Hãy đọc Chang và cộng sự để xem tác động của việc trộm cắp dữ liệu đối với giá trị thị trường của một công ty.

\section{Các vấn đề về chất lượng dữ liệu (Data Quality Problems)}

Bạn đã nghe đến từ viết tắt GIGO: Garbage In, Garbage Out (Rác vào, Rác ra). Tuy nhiên, làm thế nào chúng ta biết rằng chúng ta có rác ngay từ đầu? Mặc dù không phải là một vấn đề đạo đức trực tiếp, một số dữ liệu có chất lượng kém. Nếu bạn biết dữ liệu có vấn đề về chất lượng nhưng vẫn mô hình hóa nó một cách sâu rộng và báo cáo về các kết quả, bạn đang nhân rộng (propagate) vấn đề chất lượng đó. Sự chặt chẽ và tính hình thức của một mô hình không cải thiện chất lượng của dữ liệu vốn đã có vấn đề ngay từ đầu (xem Liu để biết thêm thông tin).

Một cách để tránh các vấn đề về chất lượng dữ liệu là yêu cầu các báo cáo tài chính đã được kiểm toán (audited financial statements). Sự giám sát của một công ty kế toán chuyên nghiệp cung cấp độ tin cậy và niềm tin vào các kết quả được báo cáo. Hãy xem xét trường hợp của công ty cho thuê xe lớn nhất thế giới, Hertz. Vài năm trước, SEC đã cáo buộc Hertz về việc báo cáo dữ liệu không chính xác (U.S. Securities and Exchange Commission, ``SEC Charges Hertz''). Bản thân Hertz đã phải trả một khoản tiền phạt. Nhưng còn các nhà đầu tư tổ chức, những người đã đưa ra các quyết định đầu tư về Hertz dựa trên thông tin này thì sao? Rõ ràng, tin tức về sự xuyên tạc đã không được biết đến vào thời điểm đó. Những người có trách nhiệm ủy thác nên hết sức cẩn thận để đánh giá chất lượng dữ liệu mà họ có trước khi họ xây dựng các mô hình và đưa ra quyết định.

\section{Thu thập dữ liệu khi chưa được phép (Collecting Data without Permission)}

Các ứng dụng điện thoại thông minh thích dữ liệu. Nhiều ứng dụng trong số đó thu thập dữ liệu của chúng ta, ngay cả khi không có sự cho phép của chúng ta. Hiện nay có một lĩnh vực tập trung vào luật quyền riêng tư dữ liệu (data privacy laws) bảo vệ người tiêu dùng, người dùng ứng dụng, v.v. khỏi các công ty thu thập và bán lại dữ liệu của họ (Klosowski). Vấn đề là dữ liệu này có thể được sử dụng theo những cách mà người dùng không bao giờ dự định.

Đạo đức tốt đòi hỏi sự cho phép rõ ràng (explicit permission) để lưu trữ dữ liệu người dùng. Sự tiết lộ này thường là một phần của thỏa thuận cấp phép (license agreement) nhưng có thể dễ dàng bị lạc trong các dòng chữ in nhỏ (fine print).

\section{Tiết lộ dữ liệu khi chưa được phép (Disclosing Data without Permission)}

Việc tiết lộ dữ liệu mà không có sự cho phép thậm chí còn tồi tệ hơn. Điều này có nghĩa là dữ liệu sử dụng của bạn được gửi cho các công ty và các bên thứ ba mà không có sự cho phép hoặc kiến thức của bạn. Hơn nữa, điều gì sẽ xảy ra nếu có một vụ vi phạm dữ liệu (data breach)? Đã có những công ty lớn mà các ngân hàng dữ liệu của họ bị hack, khiến thông tin nhạy cảm của mọi người có sẵn cho tội phạm.

Như bạn có thể thấy, dữ liệu có bộ quy tắc đạo đức riêng của nó. Với nhiều loại dữ liệu hơn có sẵn (hãy nghĩ đến tất cả các dữ liệu thay thế (alternative data) được sử dụng trong tài chính), số lượng tài khoản mà mọi người có ngày càng tăng (với ngân hàng, công ty môi giới, sàn giao dịch, cố vấn, ứng dụng, v.v.), có nhiều dữ liệu hơn bao giờ hết.

\section{Kết luận (Conclusion)}

Giống như việc chúng ta đã kết thúc khóa học FM bằng việc thảo luận về đạo đức, chúng ta kết thúc khóa học FD bằng việc thảo luận về đạo đức xoay quanh dữ liệu. Chúng tôi nhắc lại rằng danh tiếng xây dựng rất chậm ở khía cạnh tích cực nhưng lại sụp đổ nhanh chóng ở khía cạnh tiêu cực. Nhiều thập kỷ thực hành đạo đức, siêng năng có thể bị lấn át bởi một hành động phi đạo đức duy nhất. Thật không may, internet không có khả năng lãng quên. Nhờ vào các bức ảnh trên Instagram, các bài đăng trên Facebook, các blog, các video trên YouTube lan truyền (viral), và các bài báo không bao giờ biến mất, những câu chuyện về hành vi phi đạo đức không có ngày hết hạn.

Những câu chuyện chúng ta đã nghe---giao dịch nội gián (insider trading), bịa đặt email, giả mạo thị trường---đã biến những người tương đối vô danh thành các meme trên internet. Các tìm kiếm trên Google tiết lộ hàng ngàn trang về những người này, những người được thế giới biết đến vì hành động khét tiếng của họ. Chi phí cơ hội (opportunity cost) khá cao. Nhiều người sẽ tránh làm việc với hoặc thuê những người này vì thành tích đã biết là không trung thực và lừa dối. Hãy thực hành đạo đức tốt --- trong dữ liệu và trong hành động.

\section{Tài liệu tham khảo (References)}

Arliss, Will. ``Probability of the Challenger Disaster.'' Medium, 9 October 2020, \url{https://towardsdatascience.com/probability-of-the-challenger-disaster-c2a17b7984cc}.

Barret, Brian. ``An Artist Used 99 Phones to Fake a Google Maps Traffic Jam.'' \emph{Wired}, 3 Feb. 2020, \url{https://www.wired.com/story/99-phones-fake-google-maps-traffic-jam/}.

Butler, Oobah. ``I Made My Shed the Top Rated Restaurant On TripAdvisor.'' \emph{Vice}, 6 December 2017, \url{https://www.vice.com/en/article/434gqw/i-made-my-shed-the-top-rated-restaurant-on-tripadvisor}.

Chang, Kuo-Chung, et al. ``The Effect of Data Theft on a Firm's Short-Term and Long-Term Market Value.'' \emph{Mathematics}, vol. 8, no. 5, 2020, \url{https://www.mdpi.com/2227-7390/8/5/808/htm}.

D'Acunto, Francesco, and Alberto Rossi. ``Robo-advice: An Effective Tool to Reduce Inequalities?'' Brookings, 5 Oct. 2022, \url{https://www.brookings.edu/research/robo-advice-an-effective-tool-to-reduce-inequalities/}.

GaNun, Jacqueline. ``Accounting Giant Ernst Young Admits its Employees Cheated on Ethics Exams.'' NPR, 28 June 2022, \url{https://www.npr.org/2022/06/28/1108044858/accounting-giant-ernst-young-admits-its-employees-cheated-on-ethics-exams}.

Grammatikos, Theoharry, and Nikolaos I. Papanikolaou. ``Applying Benford's Law to Detect Accounting Data Manipulation in the Banking Industry.'' \emph{Journal of Financial Services Research}, vol. 59, 2021, pp. 115--142, \url{https://link.springer.com/article/10.1007/s10693-020-00334-9}.

Horn, Sebastian, et al. ``Systematic Underreporting of Public Debt Statistics: 50 Years of Evidence and Recent Progress.'' \emph{World Bank Blogs}, 3 March 2022, \url{https://blogs.worldbank.org/developmenttalk/systematic-underreporting-public-debt-statistics-50-years-evidence-and-recent}.

Klosowski, Thorin. ``The State of Consumer Data Privacy Laws in the US (And Why It Matters).'' \emph{The New York Times}, 6 Sept. 2021, \url{https://www.nytimes.com/wirecutter/blog/state-of-privacy-laws-in-us/}.

Le Marois, Oliver, and Raphael Douady. ``Return Smoothing Practices.'' \emph{The Hedge Fund Journal}, no. 30, 2007, \url{https://thehedgefundjournal.com/return-smoothing-practices/}.

Levitt, Steven D., and Stephen J. Dubner. \emph{Freakonomics: A Rogue Economist Explores the Hidden Side of Everything}. William Morrow, 2005.

Liu, Grace. ``Data Quality Problems Troubling Business and Financial Researchers: A Literature Review and Synthetic Analysis.'' \emph{Journal of Business \& Finance Librarianship}, November 2020, \url{https://digitalcommons.wcupa.edu/lib_facpub/13/}.

U.S. Securities and Exchange Commission. ``SEC Charges Investment Adviser with Failing to Disclose Conflicts Arising from Mutual Fund and Account Recommendations.'' March 2022, \url{https://www.sec.gov/litigation/litreleases/2022/lr25340.htm}.

U.S. Securities and Exchange Commission. ``SEC Charges Father-and-Son Hedge Fund Managers Who Agree to Pay \$4.8 Million to Settle Fraud Case.'' 2012, \url{https://www.sec.gov/news/press-release/2012-2012-71htm}.

U.S. Securities and Exchange Commission. ``SEC Charges Hertz with Inaccurate Financial Reporting and Other Failures.'' 1 Feb. 2019, \url{https://www.sec.gov/enforce/33-10601-s}.

U.S. Securities and Exchange Commission. ``SEC Charges Two Robo-Advisers With False Disclosures.'' 2018, \url{https://www.sec.gov/news/press-release/2018-300}.

Copyright 2024 WorldQuant University. This content is licensed solely for personal use. Redistribution or publication of this material is strictly prohibited.
